\documentclass[../sparc.tex]{subfiles}
\graphicspath{{\subfix{../images/}}}
\begin{document}

\newpage
%%%%%%%%%%%%%%%%%%%%%%%%%%%%%%%%%%%%%%%%%%%%%%%%%%%%%%%%%%%%%%%%%%%%%%%%%%%%%%%%
\section{Физика}
\index{Разработка игр!Физика}

В данном разделе мы поговорим о том, как реализовать простую физику прыжка в
нашей игре.  Для этого нам необходимо добавить в наш игровой движок понятие
гравитации и скорости.

\subsection{Гравитация и скорость}

Определим, что на позицию игрока по оси Y влияют две противоборствующие силы,
которые можно представить в виде векторов: гравитация, направленная вниз, и сила
прыжка, направленная вверх.

Гравитация определяется ускорением свободного падения.  В реальном физическом
мире, на планете Земля ускорение свободного падения составляет примерно $9.81
\mbox{м}/\mbox{с}^{2}$ или $1g$.  В нашем игровом мире это значение мы можем
подобрать по своему желанию.

Сила прыжка определяется также ускорением, направленным вверх.

Если эти две силы находятся в балансе, позиция игрока по оси Y остаётся
неизменной.  Как только баланс нарушается, то позиция начинает меняться.  Само
собой, если игрок стоит на твёрдой поверхности, это сводит на нет эффект
гравитации.  Но как только игрок оказывается в воздухе, гравитация начинает
тянуть его вниз.

Обе эти силы влияют на движение игрока по оси Y -- обозначим это понятие
скоростью, или изменением позиции в единицу времени.

Значение гравитации добавляется к значению скорости таким образом, чтобы позиция
игрока по Y изменялась, двигая его вниз.  Значение силы прыжка добавляется к
скорости, чтобы двигать игрока вверх.

\subsection{Реализация}

Зададим гравитацию (ускорение свободного падения), как вещественную константу:
\begin{minted}{cpp}
  const float GRAVITY_FORCE = 0.25;
\end{minted}

Значение константы необходимо подобрать под наши требования к игре, но пока мы
возьмём её, равной 0.25.

Далее зададим силу прыжка также вещественной константой, пусть она будет равна
0.50:
\begin{minted}{cpp}
  const float JUMP_FORCE = 0.50;
\end{minted}

Для хранения скорости по оси Y нам потребуется переменная:
\begin{minted}{cpp}
  float y_velocity = 0;
\end{minted}

Сделаем процедуру для обработки физики.  Первым делом мы будем увеличивать
скорость движения персонажа вниз, добавляя к нему наше ускорение свободного
падения -- но при условии, что игрок не стоит на платформе:

\begin{minted}{cpp}
  void handle_physics() {
    if (! is_wall(player_x, player_y + 1)) {
      y_velocity += GRAVITY_FORCE;
    }
  }
\end{minted}

Далее необходимо менять позицию игрока по оси Y на значение, которое на данный
момент хранится в переменной \mintinline{cpp}{y_velocity}.  Можно сделать
наивную реализацию этого процесса, в виде добавления значения
\mintinline{cpp}{y_velocity} к переменной \mintinline{cpp}{player_y}:

\begin{minted}{cpp}
  void handle_physics() {
    if (! is_wall(player_x, player_y + 1)) {
      y_velocity += GRAVITY_FORCE;
    }
    player_y += y_velocity;
  }
\end{minted}

Однако подобная реализация будет хорошо работать только в том случае, если
изменение позиции игрока по оси Y будет равно 1.  Если же это изменение больше
1, то тогда ирок будет проваливаться сквозь пол при падении, либо пробивать
потолок при прыжке.  Чтобы этого не происходило, необходимо придумать более
универсальное решение.

Одно из возможных решений выглядит так.  Мы можем разделить движение на 2 клети
вниз на два перемещения по 1 клетке.  Это позволит на каждом шаге проверять, не
столкнулись ли мы со стеной.

\begin{minted}{cpp}
  void handle_physics() {
    if (! is_wall(player_x, player_y + 1)) {
      y_velocity += GRAVITY_FORCE;
    }
    int delta = (y_velocity > 0) ? 1 : -1;
    for (int n = 0; n < abs(y_velocity); n++) {
      if (is_wall(player_x, player_y + delta)) {
        y_velocity = 0;
        break;
      }
      player_y += delta;
    }
  }
\end{minted}

В процедуре \mintinline{cpp}{loop} нам необходимо вызывать процедуру
\mintinline{cpp}{handle_physics}:

\begin{minted}{cpp}
void loop() {
  handle_buttons();
  handle_physics();
  map_show();
  // ...
  delay(100);
}
\end{minted}

Обработку прыжка игрока мы можем сделать следующим образом:

\begin{minted}{cpp}
  void handle_buttons() {
    // ...
    if (digitalRead(BUTTON_U) == LOW) {
      step_sound();
      if (player_y > 0) {
        y_velocity -= JUMP_FORCE;
      }
    }
    // ...
  }
\end{minted}



\end{document}
